\chapter{Le verrou : une donnée ouverte mais inexploitable}
\label{ch:verrou}

\questionchapitre{La donnée hydrologique française est publique et accessible. Pourquoi
faut-il malgré tout construire une infrastructure avant d'espérer entraîner le moindre
modèle ?}

\section{Les sources et ce qu'elles fournissent}

\textbf{Hub'Eau}~\cite{hubeau} est la plateforme d'accès aux données du système d'information
sur l'eau. Elle expose par services web les chroniques \emph{piézométriques}, c'est-à-dire les
niveaux de nappe mesurés aux ouvrages de surveillance, ainsi que les données
\emph{hydrométriques} de débit des cours d'eau, organisées en sites, stations et observations
élaborées. Le parc compte \num{23333} ouvrages piézométriques et \num{6468} stations
hydrométriques réparties sur \num{9283} sites.

\textbf{ERA5-Land}~\cite{munozsabater2021era5land}, produit par le Centre européen pour les
prévisions météorologiques à moyen terme et diffusé par le \emph{Climate Data
Store}~\cite{cds}, fournit une réanalyse climatique sur grille d'environ \ang{0.1}, continue
depuis 1950~\cite{hersbach2020era5}. C'est la source de forçage pour tout modèle
hydrologique : précipitation, température, évaporation.

\textbf{BDLISA}~\cite{bdlisa}, référentiel hydrogéologique du BRGM, décrit les entités
aquifères du territoire. C'est lui qui permet de savoir si un piézomètre observe une nappe
alluviale, un aquifère sédimentaire poreux, un système karstique ou un socle fracturé,
distinction dont le chapitre~\ref{ch:metier} montrera qu'elle commande le comportement du
signal.

Ces trois sources sont ouvertes, documentées et interrogeables. Le verrou n'est pas
l'accessibilité.

\begin{remarque}[title={Un arbitrage sur la source climatique}]
Une alternative existait. Le BRGM dispose d'une licence pour la réanalyse
SAFRAN~\cite{safran}, mieux adaptée au territoire français qu'ERA5-Land. Elle a été écartée
pour une raison qui n'est pas technique : les données SAFRAN ne sont pas redistribuables, ce
qui interdisait de les verser dans un entrepôt destiné à être partagé entre équipes et
documenté publiquement. ERA5-Land est libre, mondial et versionné. Le choix privilégie donc la
reproductibilité et le partage sur la finesse de la source, arbitrage qu'il faut connaître pour
interpréter les résultats.
\end{remarque}

\section{Ce qui manque}

\subsection{La jointure au forçage climatique n'existe pas}

Un niveau de nappe seul n'est pas modélisable. Il répond à la pluie efficace, différence entre
précipitation et évapotranspiration, avec un retard et un amortissement propres à l'aquifère.
Prédire une nappe suppose donc de disposer, pour chaque station et chaque jour, du signal
climatique correspondant.

Hub'Eau ne fournit aucune variable météorologique et ERA5 ignore l'existence des piézomètres.
Rattacher les deux exige une jointure spatiale, associant chaque station à la maille
climatique la plus proche, puis une jointure temporelle sur un pas de temps commun. Ce travail
n'est fourni par aucun des producteurs.

\subsection{L'historique n'est pas consolidé, et la volumétrie est réelle}

Les services web sont paginés et conçus pour la consultation, non pour la constitution d'un
corpus. Récupérer six décennies de chroniques sur l'ensemble du parc, puis sept décennies de
réanalyse climatique sur le territoire, représente des dizaines de gigaoctets et plusieurs
jours de collecte. Ce n'est pas une opération que l'on relance à chaque expérience.

Le cas d'ERA5 illustre l'écart entre l'accès théorique et le coût pratique. Le produit dérivé
de statistiques journalières proposé par le \emph{Climate Data Store} est calculé côté serveur,
sur une file d'attente saturée, soit environ \num{43} heures par année demandée et près de six
semaines pour reconstituer l'historique complet. La même information, recalculée localement à
partir de l'archive horaire brute, demande environ vingt-cinq minutes par année. L'équivalence
des deux voies a été vérifiée maille par maille et jour par jour, à \SI{0.0000}{\celsius} sur
les extrêmes journaliers et \SI{0.01}{\celsius} sur la moyenne.

\subsection{La donnée brute recèle des pièges silencieux}

Un exemple, parce qu'il a produit un diagnostic erroné avant d'être compris. Dans ERA5, la
précipitation et l'évaporation sont des \emph{flux cumulés} : la valeur au pas de
\SI{00}{\hour}~UTC constitue bien le total journalier. La température est une grandeur
\emph{instantanée} : prendre sa valeur au même pas comme moyenne journalière introduit un
biais froid nocturne de \SIrange{2}{4}{\celsius}. Deux variables issues du même fichier, au
même horodatage, dont l'une se lit telle quelle et l'autre doit être reconstruite à partir des
vingt-quatre pas horaires.

Rien dans la donnée ne signale cette différence. Elle se paie en résultats faux, silencieux et
plausibles.

\subsection{Aucun indicateur normalisé n'est fourni}

Un hydrogéologue ne raisonne pas sur un niveau brut en mètres, qui n'a de sens que rapporté à
l'historique de l'ouvrage. Il raisonne sur des indicateurs standardisés situant la mesure du
mois par rapport à la distribution de référence du même mois : normale, modérément sèche,
sécheresse sévère. Ces indicateurs ne sont fournis par aucune des trois sources. Ils doivent
être calculés, ce qui suppose une période de référence, un ajustement statistique et une
validation. Le chapitre~\ref{ch:indicateurs} y est consacré.

\section{Ce que l'apprentissage automatique exige en plus}

Les difficultés précédentes concernent tout usage de la donnée. L'apprentissage automatique en
ajoute qui lui sont propres.

\begin{itemize}
  \item \textbf{Une chronologie explicite.} Un modèle de séries temporelles supporte mal les
  lacunes ; encore faut-il savoir où elles sont, ce qui suppose une chronologie régulière
  plutôt qu'une liste d'observations.
  \item \textbf{Des covariables alignées.} Le forçage climatique doit être disponible sur
  exactement la même chronologie que la cible, sans décalage ni trou différentiel.
  \item \textbf{La reproductibilité.} Une expérience qui ne peut être rejouée sur les mêmes
  données n'est pas un résultat. Cela interdit l'extraction ponctuelle et impose un stockage
  versionné et interrogeable.
  \item \textbf{Le débit de lecture.} Comparer une douzaine d'architectures sur des centaines
  de stations suppose de relire les données des milliers de fois. Une chaîne d'extraction par
  service web rend l'exercice impraticable ; une base correctement indexée le rend rapide.
\end{itemize}

\section{Le coût collectif de l'absence d'infrastructure}

Le constat déterminant n'est pas technique mais organisationnel. En l'absence d'entrepôt,
chaque équipe refait l'extraction, avec ses conventions de nettoyage, sa gestion des lacunes et
son rattachement climatique. Les résultats deviennent incomparables, non par divergence de
méthode, mais par divergence de préparation des données. Le temps consacré à la collecte est
prélevé sur le temps de recherche, autant de fois qu'il y a d'équipes.

La réunion de cadrage du 13 février 2026 dresse le même constat en termes opérationnels. Elle
recense trois jeux de données candidats de provenances distinctes, un ensemble d'environ
\num{220} piézomètres réputés peu influencés par les prélèvements, un ensemble d'environ
\num{170} piézomètres annotés par des experts, et un jeu reconstitué dans le cadre de travaux
de thèse, puis conclut à la nécessité de consolider et de qualifier ces jeux, et de centraliser
données, modèles, expériences et scripts sur un espace commun~\cite{reunion_aida}.

C'est ce coût réparti que l'entrepôt décrit au chapitre~\ref{ch:entrepot} supprime.

\begin{acquisouvert}
\textbf{Acquis.} Le verrou est caractérisé. La donnée est ouverte, mais ni jointe, ni
consolidée, ni normalisée, ni prête pour l'apprentissage. La question qui organise la suite est
donc : quelle infrastructure faut-il pour que la recherche en apprentissage sur les eaux
souterraines devienne possible ?

\textbf{Ouvert.} Le périmètre retenu couvre la quantité, niveaux et débits. La qualité de
l'eau, également visée par le programme, n'a pas été traitée.
\end{acquisouvert}
